% =====================================================================
%  RESUMO EM LÍNGUA ESTRANGEIRA -- ABSTRACT (obrigatório) -- Manual, 2.1.8.2
%  É a TRADUÇÃO FIEL do resumo em português, incluindo as palavras-chave
%  (Keywords). Mesmas regras de formatação do resumo.
% =====================================================================
\begin{resumo}[Abstract]
\begin{otherlanguage*}{english}
\SingleSpacing
\noindent Dropout is a recurring problem in martial arts and combat sports
gyms, and many students quit before perceiving the physical and mental
benefits of training. This work presents RoundUp, a gamified mobile
application that aims to increase student engagement and retention and to
support instructors and gyms in managing their classes. The project was
conducted as the author's ideation process: thirteen feature ideas were
converted into requirements, built as hypotheses and, in parallel,
submitted to a questionnaire used to confirm or refute them and to reveal
new ideas. The questionnaire was answered by current and former
practitioners and, to cover missing profiles such as instructors, by 65
synthetic responses, produced by a script from 12 personas created with
the support of artificial intelligence; these responses brought new ideas
and a complementary signal, but the decision on each requirement relied on
the real data. The prototype was developed with React Native and
Firebase, has three user profiles (student, instructor and gym) and
implements all 33 functional requirements, including attendance check-in
validated by the instructor, experience points, levels, training streaks,
badges, leaderboard, missions, belt promotion records, attendance reports
and consent in compliance with the Brazilian data protection law (LGPD).
Most of the code was written by AI coding agents, while the author
focused on planning, design, and quality control of the agents and the software itself: the business rules that
protect the gamification were implemented in the database security rules
and covered by 288 automated tests, and four requirements were verified on
a device. In the real evaluation, six of the twelve evaluated ideas were
confirmed, especially the features that make progress visible, such as the
progress bar toward the next belt and the attendance history; reminders,
automatic messages, missions and the attendance check-in itself were not
confirmed, and the comments pointed to excessive notifications and data
entry as risks of abandoning the app, which the project's requirements
already addressed. The synthetic responses and the comments resulted in
nine candidate requirements for future versions. It is concluded that
gamification centered on visible progress is the most promising direction
for the application, that part of the dropout has external causes that an
app cannot solve, that development with AI agents shifts the developer's
role toward specification and verification, and that the next step is to
validate the app and the management requirements with students and
instructors in the field.

\vspace{\onelineskip}
\noindent\textbf{Keywords}: gamification; martial arts; requirements engineering; personas; AI-assisted development.
\end{otherlanguage*}
\end{resumo}